% GENERATED FILE. Edit canonical lessons, then run npm run generate:books.
% canonical-source-hash: fnv1a64:e5cda2e696428b17
% canonical-lessons: PT-C506-mergulho, PT-C506-caiaque, PT-C506-velejar, PT-C506-prancha, PT-C506-surfar, PT-R506-revisao-499-506

\chapter{Desportos de água}
\label{ch:pt-desportos-de-agua}
\hlchaptermodality{506}

\begin{chapteropening}
\textbf{By the end of this chapter:} \emph{I can talk about swimming, kayaking, sailing and surfing.}

Complete the last lesson of chapter 506: I can talk about swimming, kayaking, sailing and surfing.
\end{chapteropening}

\section[o mergulho]{o mergulho — dive; (quick) swim}

\label{lesson:PT-C506-mergulho}



\begin{quote}
Before the new one: say the Portuguese for easel, then the Portuguese for (picture) frame.
\end{quote}

\subsection*{You'll want to know: o mergulho}
\textbf{o mergulho} — \textquotedblleft{}dive; (quick) swim\textquotedblright{}.

\begin{quote}
Demos um mergulho no mar antes do almoço.
\end{quote}

\begin{quote}
\emph{We had a quick swim in the sea before lunch.}
\end{quote}

\emph{Dar um mergulho} — to have a dip.

\begin{cousinweb}
\emph{mergulho} comes from \emph{mergulhar} \textquotedblleft{}to dive\textquotedblright{}, which is usually traced to a Vulgar Latin form built on Latin \emph{mergere} \textquotedblleft{}to plunge, to dip\textquotedblright{} (Lewis \& Short).
\end{cousinweb}

\begin{grammarlens}[title={What you've built}]
One more word for telling, explaining and describing at B1.
\end{grammarlens}

\subsection*{Your turn}
\begin{itemize}
\raggedright
  \item \emph{Say it:} \emph{o mergulho}
  \item \emph{Say it:} \emph{o mergulho}, once more
  \item \emph{Say it:} a sentence of your own with \emph{o mergulho}
  \item \emph{Recall:} say the Portuguese for easel, then the Portuguese for (picture) frame, then say \emph{o mergulho} again
\end{itemize}

\begin{tcolorbox}[breakable,colback=teal!4,colframe=teal!35!black,title={Before you move on}]
What does o mergulho mean? (\textquotedblleft{}Dive; (quick) swim\textquotedblright{}.) Say it once more.
\end{tcolorbox}
\section[o caiaque]{o caiaque — kayak}

\label{lesson:PT-C506-caiaque}



\begin{quote}
Before the new one: say the Portuguese for (picture) frame, then the Portuguese for dive.
\end{quote}

\subsection*{You'll want to know: o caiaque}
\textbf{o caiaque} — \textquotedblleft{}kayak\textquotedblright{}.

\begin{quote}
No verão andamos de caiaque no rio.
\end{quote}

\begin{quote}
\emph{In summer we go kayaking on the river.}
\end{quote}

\emph{Andar de caiaque} — to go kayaking.

\begin{cousinweb}
\emph{caiaque} is usually traced to an Inuit word for a small hunting boat, probably through French or English.
\end{cousinweb}

\begin{grammarlens}[title={What you've built}]
One more word for telling, explaining and describing at B1.
\end{grammarlens}

\subsection*{Your turn}
\begin{itemize}
\raggedright
  \item \emph{Say it:} \emph{o caiaque}
  \item \emph{Say it:} \emph{o caiaque}, once more
  \item \emph{Say it:} a sentence of your own with \emph{o caiaque}
  \item \emph{Recall:} say the Portuguese for (picture) frame, then the Portuguese for dive, then say \emph{o caiaque} again
\end{itemize}

\begin{tcolorbox}[breakable,colback=teal!4,colframe=teal!35!black,title={Before you move on}]
What does o caiaque mean? (\textquotedblleft{}Kayak\textquotedblright{}.) Say it once more.
\end{tcolorbox}
\section[velejar]{velejar — to sail}

\label{lesson:PT-C506-velejar}



\begin{quote}
Before the new one: say the Portuguese for dive, then the Portuguese for kayak.
\end{quote}

\subsection*{You'll want to know: velejar}
\textbf{velejar} — \textquotedblleft{}to sail\textquotedblright{}.

\begin{quote}
Aos fins de semana, velejamos na baía.
\end{quote}

\begin{quote}
\emph{At weekends, we sail in the bay.}
\end{quote}

Built on \emph{a vela} — a sail.

\begin{cousinweb}
\emph{velejar} is built on \emph{vela} \textquotedblleft{}sail\textquotedblright{}, which is inherited from Latin \emph{vēla}, plural of \emph{vēlum} \textquotedblleft{}a sail, a cloth\textquotedblright{} (Lewis \& Short).
\end{cousinweb}

\begin{grammarlens}[title={What you've built}]
One more word for telling, explaining and describing at B1.
\end{grammarlens}

\subsection*{Your turn}
\begin{itemize}
\raggedright
  \item \emph{Say it:} \emph{velejar}
  \item \emph{Say it:} \emph{velejar}, once more
  \item \emph{Say it:} a sentence of your own with \emph{velejar}
  \item \emph{Recall:} say the Portuguese for dive, then the Portuguese for kayak, then say \emph{velejar} again
\end{itemize}

\begin{tcolorbox}[breakable,colback=teal!4,colframe=teal!35!black,title={Before you move on}]
What does velejar mean? (\textquotedblleft{}To sail\textquotedblright{}.) Say it once more.
\end{tcolorbox}
\section[a prancha]{a prancha — board (for surfing); plank}

\label{lesson:PT-C506-prancha}



\begin{quote}
Before the new one: say the Portuguese for kayak, then the Portuguese for to sail.
\end{quote}

\subsection*{You'll want to know: a prancha}
\textbf{a prancha} — \textquotedblleft{}board (for surfing); plank\textquotedblright{}.

\begin{quote}
Levou a prancha para a praia.
\end{quote}

\begin{quote}
\emph{He took his board to the beach.}
\end{quote}

\emph{Prancha de surf} — surfboard.

\begin{cousinweb}
\emph{prancha} comes from French \emph{planche} \textquotedblleft{}a board\textquotedblright{} (EtymDB), from Late Latin \emph{planca} \textquotedblleft{}a plank\textquotedblright{}.
\end{cousinweb}

\begin{grammarlens}[title={What you've built}]
One more word for telling, explaining and describing at B1.
\end{grammarlens}

\subsection*{Your turn}
\begin{itemize}
\raggedright
  \item \emph{Say it:} \emph{a prancha}
  \item \emph{Say it:} \emph{a prancha}, once more
  \item \emph{Say it:} a sentence of your own with \emph{a prancha}
  \item \emph{Recall:} say the Portuguese for kayak, then the Portuguese for to sail, then say \emph{a prancha} again
\end{itemize}

\begin{tcolorbox}[breakable,colback=teal!4,colframe=teal!35!black,title={Before you move on}]
What does a prancha mean? (\textquotedblleft{}Board (for surfing); plank\textquotedblright{}.) Say it once more.
\end{tcolorbox}
\section[surfar]{surfar — to surf}

\label{lesson:PT-C506-surfar}



\begin{quote}
Before the new one: say the Portuguese for to sail, then the Portuguese for board.
\end{quote}

\subsection*{You'll want to know: surfar}
\textbf{surfar} — \textquotedblleft{}to surf\textquotedblright{}.

\begin{quote}
Os meus primos surfam todo o ano.
\end{quote}

\begin{quote}
\emph{My cousins surf all year round.}
\end{quote}

Also \emph{fazer surf}.

\begin{cousinweb}
\emph{surfar} is built on \emph{surf}, borrowed from English.
\end{cousinweb}

\begin{grammarlens}[title={What you've built}]
One more word for telling, explaining and describing at B1.
\end{grammarlens}

\subsection*{Your turn}
\begin{itemize}
\raggedright
  \item \emph{Say it:} \emph{surfar}
  \item \emph{Say it:} \emph{surfar}, once more
  \item \emph{Say it:} a sentence of your own with \emph{surfar}
  \item \emph{Recall:} say the Portuguese for to sail, then the Portuguese for board, then say \emph{surfar} again
\end{itemize}

\begin{tcolorbox}[breakable,colback=teal!4,colframe=teal!35!black,title={Before you move on}]
What does surfar mean? (\textquotedblleft{}To surf\textquotedblright{}.) Say it once more.
\end{tcolorbox}
\section[(dialogue)]{Revisão dos capítulos 499–506}

\label{lesson:PT-R506-revisao-499-506}



\begin{quote}
Say the Portuguese for grateful and gratitude.
\end{quote}

\subsection*{Your turn}
Say these aloud:

\begin{itemize}
\raggedright
  \item grateful, gratitude, ungrateful, faithful, faithfulness
  \item sewing, to sew, knitting, to knit, to embroider
  \item gardening, to sow, to water, watering can, flowerpot
  \item to fish, fishing, bait, fish hook, to hunt
  \item chess, pack of cards, (game) board, jigsaw puzzle, crossword
  \item collection, to collect, handicrafts, clay, pottery
  \item paintbrush, watercolour, canvas, easel, (picture) frame
  \item dive, kayak, to sail, board, to surf
\end{itemize}

\begin{tcolorbox}[breakable,colback=teal!4,colframe=teal!35!black,title={Before you move on}]
Grateful? (\textbf{grato}.) Gratitude? (\textbf{a gratidão}.) Ungrateful? (\textbf{ingrato}.) Faithful? (\textbf{fiel}.) Faithfulness? (\textbf{a fidelidade}.) Sewing? (\textbf{a costura}.) To sew? (\textbf{costurar}.) Knitting? (\textbf{o tricô}.) To knit? (\textbf{tricotar}.) To embroider? (\textbf{bordar}.) Gardening? (\textbf{a jardinagem}.) To sow? (\textbf{semear}.) To water? (\textbf{regar}.) Watering can? (\textbf{o regador}.) Flowerpot? (\textbf{o vaso}.) To fish? (\textbf{pescar}.) Fishing? (\textbf{a pesca}.) Bait? (\textbf{a isca}.) Fish hook? (\textbf{o anzol}.) To hunt? (\textbf{caçar}.) Chess? (\textbf{o xadrez}.) Pack of cards? (\textbf{o baralho}.) (game) board? (\textbf{o tabuleiro}.) Jigsaw puzzle? (\textbf{o puzzle}.) Crossword? (\textbf{as palavras cruzadas}.) Collection? (\textbf{a coleção}.) To collect? (\textbf{colecionar}.) Handicrafts? (\textbf{o artesanato}.) Clay? (\textbf{o barro}.) Pottery? (\textbf{a olaria}.) Paintbrush? (\textbf{o pincel}.) Watercolour? (\textbf{a aguarela}.) Canvas? (\textbf{a tela}.) Easel? (\textbf{o cavalete}.) (picture) frame? (\textbf{a moldura}.) Dive? (\textbf{o mergulho}.) Kayak? (\textbf{o caiaque}.) To sail? (\textbf{velejar}.) Board? (\textbf{a prancha}.) To surf? (\textbf{surfar}.)
\end{tcolorbox}
